% !TEX program = xelatex
\documentclass[14pt]{extarticle}

% ============================================================
% GEOMETRY
% ============================================================
\usepackage[
  a4paper,
  left=2.5cm,
  right=1.5cm,
  top=2cm,
  bottom=2cm
]{geometry}

% ============================================================
% МОВА ТА ШРИФТИ
% ============================================================
\usepackage[ukrainian]{babel}
\usepackage{fontspec}
\setmainfont{Times New Roman}

% ============================================================
% ОСНОВНИЙ ТЕКСТ
% ============================================================
\usepackage{setspace}
\onehalfspacing
\setlength{\parindent}{1.27cm}
\setlength{\parskip}{0pt}
\usepackage{microtype}

% ============================================================
% ЗАГОЛОВКИ
% ============================================================
\usepackage{titlesec}

\titleformat{\section}[block]
  {\bfseries\fontsize{16pt}{24pt}\selectfont\centering}
  {РОЗДІЛ~\thesection.}{0.5em}{}
\titlespacing*{\section}{0pt}{20pt}{6pt}

\titleformat{\subsection}[block]
  {\bfseries\fontsize{14pt}{21pt}\selectfont}
  {\thesubsection.}{0.5em}{}
\titlespacing*{\subsection}{1.27cm}{6pt}{6pt}

\titleformat{\subsubsection}[block]
  {\bfseries\itshape\fontsize{14pt}{21pt}\selectfont}
  {}{0pt}{}
\titlespacing*{\subsubsection}{1.27cm}{6pt}{0pt}

% ============================================================
% ЗМІСТ
% ============================================================
\usepackage{tocloft}
\renewcommand{\contentsname}{\centering\bfseries\fontsize{16pt}{24pt}\selectfont ЗМІСТ}
\renewcommand{\cftsecpresnum}{РОЗДІЛ~}
\renewcommand{\cftsecaftersnum}{.}
\setlength{\cftsecnumwidth}{3.6cm}
\setlength{\cftsubsecindent}{3.6cm}
\setlength{\cftsubsecnumwidth}{1.0cm}
\renewcommand{\cftsecfont}{\bfseries}
\renewcommand{\cftsecpagefont}{\bfseries}

% ============================================================
% СПИСКИ
% ============================================================
\usepackage{enumitem}

\setlist[itemize,1]{
  label=$\bullet$,
  leftmargin=1.27cm,
  labelwidth=0.5cm,
  labelsep=0.3cm,
  itemindent=0pt,
  itemsep=3pt,
  topsep=6pt,
  parsep=0pt,
  partopsep=0pt
}

\setlist[enumerate,1]{
  label=\arabic*.,
  leftmargin=2.54cm,
  labelwidth=0.63cm,
  labelsep=0.54cm,
  itemindent=0pt,
  itemsep=0pt,
  topsep=0pt,
  parsep=0pt,
  partopsep=0pt
}

% ============================================================
% ТАБЛИЦІ
% ============================================================
\usepackage{array}
\usepackage{longtable}
\usepackage{booktabs}
\usepackage{tabularx}
\usepackage{multirow}
\usepackage[table]{xcolor}

\definecolor{riskred}{RGB}{255,102,102}
\definecolor{riskyellow}{RGB}{255,220,80}
\definecolor{riskgreen}{RGB}{102,204,102}

% ============================================================
% ДІАГРАМИ
% ============================================================
\usepackage{tikz}
\usetikzlibrary{arrows.meta,positioning,calc}
\usepackage{pgfgantt}
\usepackage{graphicx}
\usepackage{rotating}

% ============================================================
% НУМЕРАЦІЯ СТОРІНОК
% ============================================================
\usepackage{fancyhdr}
\pagestyle{fancy}
\fancyhf{}
\fancyfoot[C]{\thepage}
\renewcommand{\headrulewidth}{0pt}
\renewcommand{\footrulewidth}{0pt}

% ============================================================
% БІБЛІОГРАФІЯ
% ============================================================
\usepackage[hidelinks,breaklinks=true]{hyperref}

% ============================================================
% ДОПОМІЖНІ МАКРОСИ
% ============================================================
\newcommand{\tabhead}[1]{\textbf{#1}}

\begin{document}

\sloppy
\emergencystretch=3em
\setlength{\tabcolsep}{4pt}

% ==================== ЗМІСТ ====================
\thispagestyle{empty}
\setcounter{page}{2}
\tableofcontents
\clearpage

% ==================== ВСТУП ====================
\setcounter{page}{3}
\section*{ВСТУП}
\addcontentsline{toc}{section}{ВСТУП}

Розвиток цифрових технологій і широке впровадження методів штучного інтелекту в задачі обробки природної мови відкривають нові можливості для вирішення прикладних мовних проблем. Одним із таких завдань є нормалізація регіональних діалектів та мовних суржиків до літературної норми -- задача, яка набуває особливої актуальності в контексті масового впровадження NLP-систем в освіту, медіа та державне управління.

Українська мова характеризується значною діалектною варіативністю: гуцульський, бойківський та закарпатський діалекти мають суттєві лексичні й граматичні відмінності від літературної норми. Поряд із ними широко поширений суржик -- суміш елементів української та російської мов, що ускладнює роботу сучасних NLP-пайплайнів, налаштованих на стандартну літературну мову. Внутрішнє переміщення людей, яке різко посилилось після початку повномасштабного вторгнення, актуалізує потребу в практичних інструментах міжрегіонального мовного порозуміння.

Метою цієї розрахунково-графічної роботи є формування повної технічної та комунікаційної документації для IT-проєкту <<Система для автоматичної нормалізації суржику та діалектних форм української мови>> -- від обгрунтування доцільності й ідентифікації стейкхолдерів до аналізу ризиків, побудови мережевого графіку та закриття проєкту. Документування охоплює всі п'ять фаз управління проєктом відповідно до методології PMBOK~\cite{pmbok} і відображає реальний стан виконаних і запланованих робіт.

\clearpage

% ==================== РОЗДІЛ 1 ====================
\section{АКТУАЛЬНІСТЬ ТЕМИ ТА МЕТА РОБОТИ}

\subsection{Актуальність теми}

Проблема автоматичної нормалізації діалектних та змішаних форм мови є актуальною з кількох взаємопов'язаних причин.

\textbf{Соціальний аспект.} Внутрішнє переміщення мільйонів людей внаслідок воєнних подій призвело до зіткнення регіональних діалектів у нових для носіїв мовних середовищах. Відсутність зрозумілих цифрових інструментів перекладу і нормалізації ускладнює процеси адаптації, навчання та комунікації.

\textbf{Технологічний аспект.} Сучасні NLP-системи -- автоматичне виправлення тексту, пошукові системи, чат-боти, системи розпізнавання мовлення -- розроблені переважно для літературної норми. Введення діалектних або суржикових текстів суттєво знижує їх якість. Відсутність спеціалізованих модулів нормалізації є вузьким місцем у всьому NLP-пайплайні.

\textbf{Науковий аспект.} Для більшості діалектів української мови відсутні великі анотовані паралельні корпуси. Дослідження у цій галузі відрізняється низькоресурсністю, що ускладнює застосування стандартних методів навчання. Синтетична генерація даних за допомогою великих мовних моделей є перспективним підходом для компенсації дефіциту реальних пар.

\textbf{Конкурентний аспект.} На початок проєкту відсутні відкриті україномовні системи нормалізації, що охоплювали б декілька діалектів і суржик з використанням сучасних нейронних архітектур. Найближчі аналоги -- Vuyko Mistral та UA-GEC~\cite{uagec} -- вирішують суміжні, але не ідентичні задачі.

\subsection{Конкретизована мета та завдання проєкту}

\textbf{Мета проєкту:} розробити та експериментально оцінити систему автоматичної нормалізації українських діалектів (гуцульського, бойківського, закарпатського) та російсько-українського суржику до літературної форми на основі нейронних моделей.

Для досягнення мети необхідно виконати такі завдання:

\begin{enumerate}
  \item Провести аналіз предметної галузі та огляд існуючих підходів до нормалізації діалектних текстів.
  \item Сформувати паралельний корпус для гуцульського діалекту (ручні та синтетичні дані).
  \item Навчити baseline seq2seq модель (umt5-base) та отримати кількісні метрики якості.
  \item Розширити корпус до мульти\-діалектного через синтетичну генерацію на основі LLM.
  \item Виконати інструкційне донавчання моделі MamayLM з використанням QLoRA~\cite{qlora}.
  \item Провести порівняльне оцінювання за метриками BLEU, chrF++, TER та COMET.
  \item Розробити веб-інтерфейс на базі NiceGUI~\cite{hf_transformers} для демонстрації системи.
  \item Сформувати повну технічну та управлінську документацію проєкту.
\end{enumerate}

\clearpage

% ==================== РОЗДІЛ 2 ====================
\section{АНАЛІТИЧНИЙ ОГЛЯД НАЯВНИХ ДАНИХ}

\subsection{Стан задачі обробки діалектної мови в NLP}

Задача нормалізації діалектних і змішаних текстів є відносно новим напрямком у ширшій галузі обробки природної мови (NLP). Серед суміжних завдань виділяють: граматичну корекцію тексту (GEC), машинний переклад (MT) та лексичну нормалізацію. Різниця полягає в тому, що нормалізація діалекту вимагає системного знання морфологічних, лексичних і фонетичних відмінностей конкретного діалекту від стандартної форми мови.

Для грецької, арабської, шотландської та норвезької мов існують спеціалізовані нормалізаційні системи, проте для регіональних діалектів та суржику української мови такі системи практично відсутні. Даний проєкт є одним із перших систематичних досліджень у цій підгалузі.

Задача нормалізації формально визначається як: дано текст $x$ у діалектній формі (або суржику), знайти відповідну нормалізовану форму $y$ літературною мовою. На відміну від граматичної корекції, нормалізація діалекту потребує передусім лексичних замін, а не граматичних виправлень. Це зближує задачу з монолінгвальним машинним перекладом.

\subsection{Аналіз існуючих підходів та інструментів}

Існуючі підходи до нормалізації та суміжних задач поділяються на три класи.

\textbf{Правилові підходи (rule-based).} Використовують словники відповідностей і граматичні правила для замін. Перевагами є прозорість і можливість ручного редагування. Недоліком є обмежене покриття і нечутливість до контексту~\cite{pmbok}. Для гуцульського діалекту розроблено транскарпатський словник (7\,469 словникових статей), що є ключовим ресурсом даного проєкту.

\textbf{Статистичні та нейронні seq2seq моделі.} Застосовуються до задачі як до задачі MT: модель навчається на паралельних корпусах і відображає діалектний текст у літературний. Архітектура Transformer~\cite{vaswani} стала стандартом для подібних задач. Модель mT5~\cite{raffel} (multilingual Text-to-Text Transfer Transformer) підтримує українську мову і є основою для baseline системи проєкту (umt5-base). Helsinki-NLP пропонує попередньо навчені моделі перекладу для ряду мов.

\textbf{Інструкційне донавчання LLM.} Найсучасніший підхід -- використання великих мовних моделей із інструкційним донавчанням (instruction fine-tuning). Для мінімізації обчислювальних витрат застосовується LoRA~\cite{lora} і QLoRA~\cite{qlora}. Для україномовних задач перспективною є модель MamayLM (Gemma-3-4B-IT, INSAIT-Institute), яка пройшла передтренування та інструкційне налаштування на україномовних корпусах. Робота з нею реалізується через бібліотеки HuggingFace Transformers~\cite{hf_transformers} та PEFT.

\textbf{Суміжні результати для української мови:}
\begin{itemize}
  \item UA-GEC~\cite{uagec} -- корпус та система граматичної корекції для літературної української мови. Забезпечив корисний зразок пайплайну GEC, хоча не вирішує задачу нормалізації діалектів.
  \item Vuyko Mistral -- дослідницький проєкт перекладу між регіональними варіантами україн\-ської мови за допомогою донавченої Mistral-7B. Близький за задачею, але охоплює інші діалекти без підтримки суржику.
  \item Spivavtor -- відкритий датасет для монолінгвальних задач (перефразування, спрощення). Використаний як джерело літературних речень для синтетичної генерації діалектних пар.
\end{itemize}

За результатами аналізу у проєкті обрано гібридний підхід: umt5-base як baseline і MamayLM з QLoRA як основна система. Синтетичні дані генеруються через LLM (OpenRouter API) за правилами і словниками для кожного діалекту.

\clearpage

% ==================== РОЗДІЛ 3 ====================
\section{ОБГРУНТУВАННЯ СФЕРИ ЗАСТОСУВАННЯ\\ТА КОНЦЕПТУАЛЬНИХ ПОНЯТЬ}

\subsection{Метод «Критичних запитань»}

Метод «Критичних запитань» передбачає послідовну відповідь на ключові запитання щодо проєкту, що дозволяє визначити його доцільність, межі та зацікавлені сторони.

\subsubsection*{Питання 1. Яку проблему вирішує проєкт?}

Проєкт вирішує проблему \textbf{відсутності автоматизованого інструменту для нормалізації діалектних та суржикових форм до літературної української мови} з використанням сучасних нейромережевих підходів. Наявні засоби корекції орієнтовані на стандартні орфографічні й граматичні помилки, а не на семантичну нормалізацію діалектних конструкцій.

\subsubsection*{Питання 2. Яку потребу задовольняє проєкт?}

Проєкт задовольняє потреби трьох груп: (1)~кінцеві користувачі (студенти, мігранти, носії діалектів) потребують доступного інструменту нормалізації коротких текстів; (2)~NLP-розробники потребують модуля нормалізації як компоненту пайплайнів; (3)~дослідники потребують відтворюваного baseline-рішення для задачі діалектної нормалізації.

\subsubsection*{Питання 3. Яка сфера застосування?}

Сфера застосування: нормалізація коротких текстів (повідомлення, коментарі, чати) з гуцульського, бойківського, закарпатського діалектів та суржику до літературної форми. Система надає веб-інтерфейс для вибору типу вхідного тексту та може використовуватись як окремий інструмент або компонент NLP-системи.

\subsubsection*{Питання 4. Хто є зацікавленими сторонами?}

Носії діалектів і суржикомовці, NLP-дослідники, розробники програмного забезпечення, науковий керівник, кафедра. Повний перелік наведено у Розділі~4.

\subsubsection*{Питання 5. Що поза межами проєкту?}

Поза межами проєкту: переклад з інших мов (російської, польської тощо); нормалізація довгих документів; розпізнавання мовлення (ASR); комерційне масштабування; покриття всіх діалектів України.

\subsubsection*{Питання 6. Трикутник конкурентних обмежень}

\textbf{Час:} один навчальний семестр (лютий-червень 2026). \textbf{Ресурси:} обмежений GPU (QLoRA знижує вимоги до пам'яті); дані синтетично розширені через Mistral Small 24B. \textbf{Якість:} BLEU, chrF++, TER; umt5-base BLEU~=~74.68 на гуцульському, 89.42 на суржику. Пріоритет: якість $>$ терміни $>$ ресурси.

\subsection{Метод SMART}

Метод SMART перевіряє мету проєкту за п'ятьма критеріями.

\textbf{S -- Specific (Конкретність).} Мета сформульована конкретно: розробити систему нормалізації для 4 різновидів (гуцульський, бойківський, закарпатський діалекти та суржик) з веб-інтерфейсом, на основі seq2seq baseline і MamayLM з QLoRA. Визначені вхід, вихід, підхід та очікуваний результат.

\textbf{M -- Measurable (Вимірюваність).} Результати вимірюються: BLEU (umt5-base досяг 74.68 на гуцульському та 89.42 на суржику); chrF++ (символьна F-міра, краща для морфологічно багатих мов); TER (кількість редагувань). Контрольні набори сформовані з вручну верифікованих пар (50/діалект).

\textbf{A -- Assignable (Виконавець).} Студент групи ШІ-42 Коваль Денис Петрович суміщає ролі менеджера, розробника та аналітика. Науковий керівник Сивоконь Олексій Олегович забезпечує методичний нагляд і рецензування. Виконавець має компетенції у Python, PyTorch та NLP.

\textbf{R -- Realistic (Реалістичність).} Обидві моделі навчено і оцінено: umt5-base досяг BLEU~=~74.68 на гуцульському; MamayLM~QLoRA -- BLEU~=~51.81 на бойківському. Синтетична генерація через Mistral Small 24B забезпечила 125\,000 навчальних пар. QLoRA дозволяє донавчати 4B-модель на GPU з~$\geq$12~GB VRAM. Результати підтвердили реалістичність підходу.

\textbf{T -- Time-related (Часові рамки).} Семестровий проєкт із чітким розкладом: аналіз літератури (3 тиж.) $\to$ підготовка даних (6 тиж.) $\to$ донавчання (4 тиж.) $\to$ оцінювання (2 тиж.) $\to$ документування (4 тиж.). Загальний критичний шлях -- 19 тижнів.

\subsection{Концептуальні поняття проєкту}

\textbf{Визначення проєкту.} Проєкт -- тимчасове зусилля зі створення унікального програмного продукту: нейромережевої системи, що перетворює тексти, написані українськими діалектами та суржиком, на літературну форму. Реалізується у межах бакалаврської кваліфікаційної роботи за спеціальністю 122 <<Комп'ютерні науки>> у НУ <<Львівська Політехніка>>.

\textbf{Команда проєкту:}
\begin{itemize}
  \item \textbf{Виконавець / Менеджер:} Коваль Денис Петрович, ШІ-42. Планування, реалізація, контроль якості, документування.
  \item \textbf{Науковий керівник (Спонсор):} Сивоконь Олексій Олегович, асистент кафедри СШІ. Методичне спрямування, рецензування, погодження вимог.
\end{itemize}

\textbf{Ромб конкурентних обмежень:}
\begin{itemize}
  \item \textbf{Кошти:} нульовий фінансовий бюджет; використовуються відкриті моделі та API.
  \item \textbf{Час:} один семестр, фіксована дата захисту БКР.
  \item \textbf{Якість:} BLEU, chrF++, TER, COMET на ручних контрольних наборах.
  \item \textbf{Результат:} працюча система з веб-інтерфейсом, відкритий репозиторій, таблиця метрик.
\end{itemize}

\subsection{Складність проєкту з точки зору управління}

Відповідно до матриці типів проєктів, проєкт класифіковано як \textbf{еволюційний із середнім рівнем виклику}.

\textbf{Еволюційні вимоги.} Задача нормалізації діалектів є недостатньо формалізованою. Точну якість для кожного діалекту неможливо передбачити на початку проєкту. Вимоги уточнюються по ходу роботи: після baseline-етапу стало зрозуміло, що мульти\-діалектний корпус потребує ретельнішого стратифікованого балансування.

\textbf{Рівень виклику -- середній.} Більшість ключових компетенцій (Python, PyTorch, HuggingFace, NLP) наявні у виконавця. Водночас донавчання великих мовних моделей і специфіка низькоресурсних діалектних корпусів виходять за межі стандартного університетського курсу.

\textbf{Критичність -- менеджментний рівень:} проєкт є ключовим для досягнення академічної мети (захист БКР), але не є критичним для виживання організації. Водночас результати мають практичну цінність для NLP-спільноти та потенційних інтеграторів.

\clearpage

% ==================== РОЗДІЛ 4 ====================
\section{ЗАЦІКАВЛЕНІ СТОРОНИ ТА ФОРМУВАННЯ ВИМОГ}

\subsection{Повний перелік зацікавлених сторін}

Стейкхолдером є будь-хто, хто впливає на проєкт або на кого впливає проєкт -- позитивно чи негативно. Нижче наведено повний перелік зацікавлених сторін із оцінкою потужності (Power) та зацікавленості (Interest).

\vspace{6pt}
\noindent\textbf{Таблиця 1 --- Зацікавлені сторони проєкту}

\begin{longtable}{|p{0.04\textwidth}|p{0.25\textwidth}|p{0.17\textwidth}|p{0.11\textwidth}|p{0.11\textwidth}|p{0.14\textwidth}|}
\hline
\tabhead{ID} & \tabhead{Стейк\-холдер} & \tabhead{Роль у проєкті} & \tabhead{Power} & \tabhead{Inter.} & \tabhead{За / Проти} \\
\hline
\endfirsthead
\hline
\tabhead{ID} & \tabhead{Стейк\-холдер} & \tabhead{Роль у проєкті} & \tabhead{Power} & \tabhead{Inter.} & \tabhead{За / Проти} \\
\hline
\endhead
S1 & Коваль Денис Петрович & Виконавець, менеджер & Висока & Висока & За \\
\hline
S2 & Сивоконь Олексій Олегович & Науковий керівник & Висока & Висока & За \\
\hline
S3 & Кафедра СШІ, НУ <<ЛП>> & Інституційний замовник & Висока & Середня & За \\
\hline
S4 & Команда MamayLM (INSAIT) & Розробники базової моделі & Низька & Середня & За \\
\hline
S5 & Носії діалектів (гуцульського, бойківського, закарпатського) & Кінцеві кори\-стувачі & Низька & Висока & За \\
\hline
S6 & Внутрішні мігранти (суржико\-мовці) & Кінцеві кори\-стувачі & Низька & Висока & За \\
\hline
S7 & Ukrainian NLP-спільнота & Академічна аудиторія & Середня & Висока & За \\
\hline
S8 & EdTech-платформи (Prometheus, МОН) & Потенційні інтегратори & Середня & Середня & Нейтрально \\
\hline
S9 & Мовні активісти / ТМШ & Промоутери норми & Низька & Висока & За \\
\hline
S10 & NLP-розробники у tech-ком\-паніях & Спо\-жи\-вачі API/ком\-по\-нен\-та & Середня & Середня & За \\
\hline
\end{longtable}

\subsection{Сітка оцінки зацікавлених сторін (Power/Interest)}

Сітка розміщує стейкхолдерів у чотирьох квадрантах за осями «Потужність» і «Зацікавленість» та визначає стратегію взаємодії.

\begin{center}
\begin{tabular}{|p{0.01\textwidth}|p{0.41\textwidth}|p{0.41\textwidth}|}
\hline
 & \multicolumn{1}{c|}{\textbf{Interest: Низька}} & \multicolumn{1}{c|}{\textbf{Interest: Висока}} \\
\hline
\rotatebox[origin=c]{90}{\textbf{Power: Висока}} &
\textbf{Квадрант II - Задовольняти} \newline S3 Кафедра СШІ / НУ «ЛП» &
\textbf{Квадрант I - Ключові гравці} \newline S1 Коваль Денис \newline S2 Сивоконь О.О. \\
\hline
\rotatebox[origin=c]{90}{\textbf{Power: Низька}} &
\textbf{Квадрант III - Моніторинг} \newline S4 Команда MamayLM &
\textbf{Квадрант IV - Інформувати} \newline S5 Носії діалектів \newline S6 Мігранти/пересе\-ленці \newline S7 Ukrainian NLP-спільнота \newline S9 Мовні активісти / ТМШ \\
\hline
\end{tabular}
\end{center}

\medskip
Стейкхолдери S8 та S10 (середній рівень за обома осями) залучаються ситуативно.

\subsection{Вимоги стейкхолдерів та фази виконання}

\vspace{6pt}
\noindent\textbf{Таблиця 2 --- Вимоги стейкхолдерів та фази виконання}

\begin{longtable}{|p{0.03\textwidth}|p{0.24\textwidth}|p{0.38\textwidth}|p{0.17\textwidth}|}
\hline
\tabhead{ID} & \tabhead{Стейкхолдер} & \tabhead{Вимоги та очікування} & \tabhead{Фаза} \\
\hline
\endfirsthead
\hline
\tabhead{ID} & \tabhead{Стейкхолдер} & \tabhead{Вимоги та очікування} & \tabhead{Фаза} \\
\hline
\endhead
S1 & Виконавець & Чітка постановка задачі, доступ до ресурсів, відтворювані результати & Всі \\
\hline
S2 & Науковий керівник & Дотримання академічних стандартів, новизна, відповідність темі БКР & Ініц., Контр., Закр. \\
\hline
S3 & Кафедра СШІ & Вчасна здача роботи, відповідність структурі БКР & Ініц., Закр. \\
\hline
S4 & Команда MamayLM & Коректне цитування моделі, зворотній зв'язок & Вик., Закр. \\
\hline
S5 & Носії діалектів & Коректна нормалізація без спотворення змісту & Вик., Контр. \\
\hline
S6 & Мігранти / переселенці & Простий веб-інтерфейс, підтримка суржику & Вик., Контр. \\
\hline
S7 & Ukrainian NLP-спільнота & Відкритий код і датасети, відтворювані метрики & Вик., Закр. \\
\hline
S8 & EdTech-платформи & Можливість інтеграції через API, документація & Закр. \\
\hline
S9 & Мовні активісти & Нормалізація до літературної норми (не «русифікація») & Вик., Контр. \\
\hline
S10 & NLP-розробники & Документований API, відкрита модель на HuggingFace & Закр. \\
\hline
\end{longtable}

\subsection{Матриця розподілу відповідальності (RAM)}

Умовні позначення: \textbf{A} -- відповідальний (Accountable); \textbf{P} -- учасник (Participant); \textbf{O} -- необхідна думка (Opinion required); порожньо -- не залучений.

\vspace{6pt}
\noindent\textbf{Таблиця 3 --- Матриця RAM}

\begin{center}
\begin{longtable}{|p{0.26\textwidth}|c|c|c|c|c|c|c|c|c|c|}
\hline
\tabhead{Завдання / Етап} & \tabhead{S1} & \tabhead{S2} & \tabhead{S3} & \tabhead{S4} & \tabhead{S5} & \tabhead{S6} & \tabhead{S7} & \tabhead{S8} & \tabhead{S9} & \tabhead{S10} \\
\hline
\endfirsthead
\hline
\tabhead{Завдання / Етап} & \tabhead{S1} & \tabhead{S2} & \tabhead{S3} & \tabhead{S4} & \tabhead{S5} & \tabhead{S6} & \tabhead{S7} & \tabhead{S8} & \tabhead{S9} & \tabhead{S10} \\
\hline
\endhead
Визначення теми та мети & A & P & O & & & & & & & \\
\hline
Огляд літератури та джерел & A & O & & & & & P & & & \\
\hline
Аналіз наборів даних & A & O & & O & & & O & & & \\
\hline
Підготовка корпусу (гуцульський) & A & O & & & O & & & & & \\
\hline
Синтетичне розширення даних & A & O & & O & O & O & & & O & \\
\hline
Навчання baseline seq2seq & A & O & & & & & & & & \\
\hline
Форм. мульти\-діалект. корпусу & A & O & & & O & O & O & & O & \\
\hline
Дон. MamayLM (QLoRA) & A & O & & P & & & & & & \\
\hline
Оцінювання якості & A & P & & & O & O & O & & O & \\
\hline
Розробка NiceGUI-інтерфейсу & A & O & & & & P & & O & & O \\
\hline
Написання тексту БКР & A & P & O & & & & & & & \\
\hline
Публікація коду та моделі & A & O & & O & & & P & O & & O \\
\hline
Захист БКР & A & P & A & & & & & & & \\
\hline
\end{longtable}
\end{center}

\clearpage

% ==================== РОЗДІЛ 5 ====================
\section{ПЛАНУВАННЯ ТА КОНТРОЛЬНІ СПИСКИ ПРОЄКТУ}

\subsection{Контрольні списки фаз PMBOK}

Проєкт проходить п'ять фаз управління відповідно до методології PMBOK~\cite{pmbok}. Для кожної фази сформовано технічний контрольний список (checklist) із очікуваними результатами.

\vspace{4pt}
\noindent\textbf{Таблиця 4 --- Фаза 1: Ініціювання}

\begin{longtable}{|p{0.53\textwidth}|p{0.33\textwidth}|}
\hline
\tabhead{Checklist-елемент} & \tabhead{Результат} \\
\hline
\endfirsthead
\hline
\tabhead{Checklist-елемент} & \tabhead{Результат} \\
\hline
\endhead
Визначення проблеми & Опис проблеми та актуальності \\
\hline
Визначення сфери застосування & Документ сфери застосування \\
\hline
Формулювання мети за SMART & Конкретна, вимірювана мета \\
\hline
Визначення зацікавлених сторін & Перелік 10 стейкхолдерів \\
\hline
Узгодження теми з науковим керівником & Затверджена тема БКР \\
\hline
Попередній огляд доступних ресурсів & Перелік ключових джерел \\
\hline
\end{longtable}

\vspace{4pt}
\noindent\textbf{Таблиця 5 --- Фаза 2: Планування}

\begin{longtable}{|p{0.53\textwidth}|p{0.33\textwidth}|}
\hline
\tabhead{Checklist-елемент} & \tabhead{Результат} \\
\hline
\endfirsthead
\hline
\tabhead{Checklist-елемент} & \tabhead{Результат} \\
\hline
\endhead
Огляд літератури (NMT, GEC, діалектна нормалізація) & Аналітичний огляд джерел \\
\hline
Аналіз наборів даних & Опис джерел даних та їх ролей \\
\hline
Вибір архітектур (umt5-base, MamayLM з QLoRA) & Обгрунтований вибір моделей \\
\hline
Визначення метрик якості & Набір метрик з цільовими значеннями \\
\hline
Визначення технологічного стеку & Перелік інструментів \\
\hline
Формування розкладу етапів & План-графік робіт \\
\hline
Ідентифікація ризиків & Реєстр ризиків \\
\hline
Визначення вимог до ресурсів & Специфікація (CPU для baseline, GPU $\geq$12~GB для LLM) \\
\hline
\end{longtable}

\vspace{4pt}
\noindent\textbf{Таблиця 6 --- Фаза 3: Виконання}

\begin{longtable}{|p{0.53\textwidth}|p{0.33\textwidth}|}
\hline
\tabhead{Checklist-елемент} & \tabhead{Результат} \\
\hline
\endfirsthead
\hline
\tabhead{Checklist-елемент} & \tabhead{Результат} \\
\hline
\endhead
Збір та підготовка даних (гуцульський) & Паралельний корпус + словник \\
\hline
Навчання baseline seq2seq (umt5-base) & Модель, BLEU~$\approx$~51 \\
\hline
Аналіз помилок baseline & Звіт про типові помилки \\
\hline
Синтетична генерація пар для 3 діалектів + суржику & Мульти\-діалектний корпус \\
\hline
Донавчання MamayLM з QLoRA & Донавчена модель \\
\hline
Оцінювання якості на контрольних наборах & Таблиця метрик \\
\hline
Розробка NiceGUI-інтерфейсу & Демо-інтерфейс \\
\hline
\end{longtable}

\vspace{4pt}
\noindent\textbf{Таблиця 7 --- Фаза 4: Контроль}

\begin{longtable}{|p{0.53\textwidth}|p{0.33\textwidth}|}
\hline
\tabhead{Checklist-елемент} & \tabhead{Результат} \\
\hline
\endfirsthead
\hline
\tabhead{Checklist-елемент} & \tabhead{Результат} \\
\hline
\endhead
Відстеження відповідності розкладу & Звіт про дотримання термінів \\
\hline
Моніторинг метрик під час навчання (TensorBoard) & Графіки loss і BLEU по епохах \\
\hline
Перевірка якості синтетичних даних & Лог відфільтрованих прикладів \\
\hline
Контроль розбиття train/val/test & Підтвердження коректності \\
\hline
Відстеження змін у вимогах & Журнал змін \\
\hline
Регулярні консультації з науковим керівником & Протоколи обговорень \\
\hline
\end{longtable}

\vspace{4pt}
\noindent\textbf{Таблиця 8 --- Фаза 5: Закриття}

\begin{longtable}{|p{0.53\textwidth}|p{0.33\textwidth}|}
\hline
\tabhead{Checklist-елемент} & \tabhead{Результат} \\
\hline
\endfirsthead
\hline
\tabhead{Checklist-елемент} & \tabhead{Результат} \\
\hline
\endhead
Оформлення тексту БКР & Готовий документ \\
\hline
Фіксація фінальних метрик по всіх варіантах & Таблиця результатів \\
\hline
Публікація коду та моделі на GitHub / HuggingFace & Відкритий репозиторій \\
\hline
Підготовка презентації до захисту & Слайди \\
\hline
Фіксація досвіду (lessons learned) & Розділ висновків у БКР \\
\hline
Захист бакалаврської роботи & Оцінка \\
\hline
\end{longtable}

\subsection{Структура розбиття робіт (WBS)}

\noindent\textbf{0.\ Система автоматичної нормалізації суржику та діалектних форм}

\vspace{4pt}
\begin{enumerate}[label=\textbf{\arabic*.}, leftmargin=1.5cm, itemsep=3pt, topsep=4pt, parsep=0pt]
  \item \textbf{Управління проєктом}
  \begin{enumerate}[label=1.\arabic*., leftmargin=1.2cm, itemsep=1pt, topsep=2pt, parsep=0pt, partopsep=0pt]
    \item Координація та звітність
    \item Управління ризиками
  \end{enumerate}

  \item \textbf{Дослідження та аналіз}
  \begin{enumerate}[label=2.\arabic*., leftmargin=1.2cm, itemsep=1pt, topsep=2pt, parsep=0pt, partopsep=0pt]
    \item Огляд академічних джерел та NLP-корпусів
    \item Аналіз наявних підходів до нормалізації
  \end{enumerate}

  \item \textbf{Підготовка даних}
  \begin{enumerate}[label=3.\arabic*., leftmargin=1.2cm, itemsep=1pt, topsep=2pt, parsep=0pt, partopsep=0pt]
    \item Збір діалектних зразків та суржик-текстів
    \item Розробка правил та словників нормалізації
    \item Синтетична генерація через LLM та анотація
  \end{enumerate}

  \item \textbf{Розробка моделей}
  \begin{enumerate}[label=4.\arabic*., leftmargin=1.2cm, itemsep=1pt, topsep=2pt, parsep=0pt, partopsep=0pt]
    \item Seq2seq baseline модель (umt5-base)
    \item Інструкційне донавчання MamayLM (QLoRA)
  \end{enumerate}

  \item \textbf{Оцінка та інтеграція}
  \begin{enumerate}[label=5.\arabic*., leftmargin=1.2cm, itemsep=1pt, topsep=2pt, parsep=0pt, partopsep=0pt]
    \item Автоматична оцінка (BLEU, chrF++, TER, COMET)
    \item Ручна оцінка вибірки (50 речень)
    \item Розробка NiceGUI-інтерфейсу
  \end{enumerate}

  \item \textbf{Документація та захист}
  \begin{enumerate}[label=6.\arabic*., leftmargin=1.2cm, itemsep=1pt, topsep=2pt, parsep=0pt, partopsep=0pt]
    \item Бакалаврська кваліфікаційна робота
    \item Технічна документація (README, CHANGELOG)
    \item Захист проєкту
  \end{enumerate}
\end{enumerate}

\subsection{Мережева діаграма та розрахунок критичного шляху}

Мережева діаграма побудована методом критичного шляху (CPM). Перелік операцій наведено в Таблиці~9.

\vspace{6pt}
\noindent\textbf{Таблиця 9 --- Перелік операцій проєкту}

\begin{longtable}{|>{\centering\arraybackslash}p{0.04\textwidth}|>{\raggedright\arraybackslash}p{0.34\textwidth}|>{\centering\arraybackslash}p{0.10\textwidth}|>{\centering\arraybackslash}p{0.13\textwidth}|>{\raggedright\arraybackslash}p{0.22\textwidth}|}
\hline
\tabhead{ID} & \tabhead{Операція} & \tabhead{Трив., тижні} & \tabhead{Залежить від} & \tabhead{Відповідальний} \\
\hline
\endfirsthead
\hline
\tabhead{ID} & \tabhead{Операція} & \tabhead{Трив., тижні} & \tabhead{Залежить від} & \tabhead{Відповідальний} \\
\hline
\endhead
A & Аналіз літератури та предметної галузі & 3 & --- & Коваль Д.П. \\
\hline
B & Підготовка та анотація датасету & 6 & A & Коваль Д.П. \\
\hline
C & Seq2seq базова модель & 3 & B & Коваль Д.П. \\
\hline
D & Донавчання MamayLM & 4 & B & Коваль Д.П. \\
\hline
E & Оцінка та інтеграція & 2 & C, D & Коваль Д.П., Сивоконь О.О. \\
\hline
F & Написання дипломної роботи & 4 & E & Коваль Д.П., Сивоконь О.О. \\
\hline
\end{longtable}

\vspace{8pt}
\noindent\textbf{Рисунок 1 --- Мережева діаграма проєкту (метод CPM)}

\vspace{8pt}
\begin{center}
\resizebox{0.98\textwidth}{!}{%
\begin{tikzpicture}[
  taskbox/.style={
    draw=gray!70, rectangle, rounded corners=3pt,
    minimum width=3.0cm, minimum height=2.8cm,
    align=center, thick
  },
  critbox/.style={
    draw=red!80!black, rectangle, rounded corners=3pt,
    minimum width=3.0cm, minimum height=2.8cm,
    align=center, thick, line width=2pt,
    fill=red!8
  },
  arr/.style={-{Stealth[length=9pt]}, thick, gray!60},
  critarr/.style={-{Stealth[length=9pt]}, thick, red!80!black, line width=2pt}
]

\node[critbox] (A) at (0, 0) {
  \large\textbf{A}\\[3pt]
  \normalsize Аналіз\\[-1pt]
  \normalsize літератури\\[5pt]
  \small ES=0 \quad EF=3\\
  \small LS=0 \quad LF=3\\[2pt]
  \small TF=0
};

\node[critbox] (B) at (5, 0) {
  \large\textbf{B}\\[3pt]
  \normalsize Підготовка\\[-1pt]
  \normalsize даних\\[5pt]
  \small ES=3 \quad EF=9\\
  \small LS=3 \quad LF=9\\[2pt]
  \small TF=0
};

\node[taskbox] (C) at (10, 2.8) {
  \large\textbf{C}\\[3pt]
  \normalsize Seq2seq\\[-1pt]
  \normalsize модель\\[5pt]
  \small ES=9 \quad EF=12\\
  \small LS=10 \; LF=13\\[2pt]
  \small TF=1
};

\node[critbox] (D) at (10, -2.8) {
  \large\textbf{D}\\[3pt]
  \normalsize Дон.\! MamayLM\\[-1pt]
  \normalsize (QLoRA)\\[5pt]
  \small ES=9 \quad EF=13\\
  \small LS=9 \quad LF=13\\[2pt]
  \small TF=0
};

\node[critbox] (E) at (15.5, 0) {
  \large\textbf{E}\\[3pt]
  \normalsize Оцінка та\\[-1pt]
  \normalsize інтеграція\\[5pt]
  \small ES=13 \; EF=15\\
  \small LS=13 \; LF=15\\[2pt]
  \small TF=0
};

\node[critbox] (F) at (20.5, 0) {
  \large\textbf{F}\\[3pt]
  \normalsize Дипломна\\[-1pt]
  \normalsize робота\\[5pt]
  \small ES=15 \; EF=19\\
  \small LS=15 \; LF=19\\[2pt]
  \small TF=0
};

\draw[critarr] (A.east) -- (B.west);
\draw[arr] (B.east) -- (C.west);
\draw[critarr] (B.east) -- (D.west);
\draw[arr] (C.east) -- (E.west);
\draw[critarr] (D.east) -- (E.west);
\draw[critarr] (E.east) -- (F.west);

\end{tikzpicture}%
}
\end{center}

\vspace{4pt}
\noindent{\small\textit{Умовні позначення: червона рамка - критичні операції (TF=0); сіра рамка - некритичні (TF$>$0).}}

\vspace{6pt}
\noindent\textbf{Таблиця 10 --- Розрахунок раннього та пізнього початку і закінчення}

\begin{longtable}{|>{\centering\arraybackslash}p{0.04\textwidth}|>{\centering\arraybackslash}p{0.09\textwidth}|>{\centering\arraybackslash}p{0.06\textwidth}|>{\centering\arraybackslash}p{0.06\textwidth}|>{\centering\arraybackslash}p{0.06\textwidth}|>{\centering\arraybackslash}p{0.06\textwidth}|>{\centering\arraybackslash}p{0.06\textwidth}|>{\centering\arraybackslash}p{0.10\textwidth}|>{\centering\arraybackslash}p{0.06\textwidth}|}
\hline
\tabhead{ID} & \tabhead{Трив.} & \tabhead{ES} & \tabhead{EF} & \tabhead{LS} & \tabhead{LF} & \tabhead{TF} & \tabhead{Крит.?} & \tabhead{FF} \\
\hline
\endfirsthead
\hline
\tabhead{ID} & \tabhead{Трив.} & \tabhead{ES} & \tabhead{EF} & \tabhead{LS} & \tabhead{LF} & \tabhead{TF} & \tabhead{Крит.?} & \tabhead{FF} \\
\hline
\endhead
A & 3 & 0 & 3 & 0 & 3 & 0 & Так & 0 \\
\hline
B & 6 & 3 & 9 & 3 & 9 & 0 & Так & 0 \\
\hline
C & 3 & 9 & 12 & 10 & 13 & 1 & Ні & 1 \\
\hline
D & 4 & 9 & 13 & 9 & 13 & 0 & Так & 0 \\
\hline
E & 2 & 13 & 15 & 13 & 15 & 0 & Так & 0 \\
\hline
F & 4 & 15 & 19 & 15 & 19 & 0 & Так & 0 \\
\hline
\end{longtable}

\noindent де: \textbf{ES} - ранній початок; \textbf{EF} - раннє закінчення; \textbf{LS} - пізній початок; \textbf{LF} - пізнє закінчення; \textbf{TF} - загальний резерв часу; \textbf{FF} - вільний резерв часу.

\medskip
\noindent Загальна тривалість проєкту: \textbf{19 тижнів}. Критичний шлях: $\mathbf{A \to B \to D \to E \to F}$.

\subsection{Матриця комунікацій проєкту}

\vspace{6pt}
\noindent\textbf{Таблиця 11 --- Матриця комунікацій}

\begin{longtable}{|>{\raggedright\arraybackslash}p{0.16\textwidth}|>{\raggedright\arraybackslash}p{0.18\textwidth}|>{\raggedright\arraybackslash}p{0.14\textwidth}|>{\raggedright\arraybackslash}p{0.13\textwidth}|>{\raggedright\arraybackslash}p{0.14\textwidth}|>{\raggedright\arraybackslash}p{0.10\textwidth}|}
\hline
\tabhead{Аудиторія} & \tabhead{Повідомлення} & \tabhead{Тип} & \tabhead{Частота} & \tabhead{Ко\-му\-ні\-ка\-тор} & \tabhead{Від\-пов.} \\
\hline
\endfirsthead
\hline
\tabhead{Аудиторія} & \tabhead{Повідомлення} & \tabhead{Тип} & \tabhead{Частота} & \tabhead{Ко\-му\-ні\-ка\-тор} & \tabhead{Від\-пов.} \\
\hline
\endhead
Нау\-ко\-вий керівник (S2) & Звіти про прогрес та технічні питання & Зустріч / email & Що\-тиж\-не\-во & Коваль Д.П. & Звіт\-ність \\
\hline
Ка\-фед\-ра СШІ (S3) & Проміжні результати & Презентація & Що\-мі\-ся\-ця & Коваль Д.П. & Від\-по\-від\-ність \\
\hline
NLP-спільнота (S7) & Публікація моделей та датасетів & GitHub / звіт & За потреби & Коваль Д.П. & Від\-кри\-тість \\
\hline
Носії діалектів (S5) & Збір зразків та зворотній зв'язок & Анкета & 1 раз & Коваль Д.П. & Якість даних \\
\hline
Команда MamayLM (S4) & Технічні консультації & Email / API & За потреби & Коваль Д.П. & Ін\-тег\-ра\-ція \\
\hline
EdTech-платформи (S8) & Демонстрація системи & Вебінар & 1 раз & Коваль Д.П.,\newline Сивоконь О.О. & Попу\-ляр. \\
\hline
\end{longtable}

\subsection{Діаграма Ганта}

\noindent\textbf{Рисунок 2 --- Діаграма Ганта проєкту (тижні 1-19)}

\vspace{8pt}
\begin{center}
\resizebox{\textwidth}{!}{%
\begin{ganttchart}[
  hgrid,
  vgrid={*{1}{draw=gray!30}},
  x unit=0.70cm,
  y unit title=0.80cm,
  y unit chart=0.90cm,
  title/.style={fill=gray!25, draw=black, font=\normalsize},
  title label font=\normalsize,
  bar label font=\normalsize,
  bar/.style={fill=blue!55, draw=blue!75!black},
  bar height=0.55,
  group right shift=0,
  group top shift=0.7,
  group height=0.3
]{1}{19}
  \gantttitle{Тижні проєкту}{19} \\
  \gantttitlelist{1,...,19}{1} \\
  \ganttbar[bar/.style={fill=red!60, draw=red!80!black}]
    {A: Аналіз літератури}{1}{3} \\
  \ganttbar[bar/.style={fill=red!60, draw=red!80!black}]
    {B: Підготовка даних}{4}{9} \\
  \ganttbar[bar/.style={fill=blue!55, draw=blue!75!black}]
    {C: Seq2seq модель}{10}{12} \\
  \ganttbar[bar/.style={fill=red!60, draw=red!80!black}]
    {D: Дон. MamayLM}{10}{13} \\
  \ganttbar[bar/.style={fill=red!60, draw=red!80!black}]
    {E: Оцінка та інтеграція}{14}{15} \\
  \ganttbar[bar/.style={fill=red!60, draw=red!80!black}]
    {F: Дипломна робота}{16}{19} \\
\end{ganttchart}%
}
\end{center}

\vspace{4pt}
\noindent{\small\textit{Умовні позначення: \textcolor{red!70!black}{\rule{1em}{0.8em}} - критичні операції ($A\to B\to D\to E\to F$); \textcolor{blue!70!black}{\rule{1em}{0.8em}} - некритичні (TF$>$0).}}

\clearpage

% ==================== РОЗДІЛ 6 ====================
\section{АНАЛІЗ РИЗИКІВ ПРОЄКТУ}

\subsection{Ідентифікація ризиків}

Управління ризиками є одним із найбільш критичних факторів успіху проєкту. Для ідентифікації ризиків проєкту застосовано такі методи: огляд аналогічних NLP-проєктів та академічних публікацій; консультація з науковим керівником щодо технічних обмежень; аналіз блок-схеми пайплайну генерації даних і тренування моделі; проведення «негативного мозкового штурму» для кожного з п'яти етапів проєкту; SWOT-аналіз специфіки роботи з мало\-ресурсними діалектними корпусами.

За результатами ідентифікації визначено 12 ризиків, що охоплюють технічні, ресурсні, організаційні та якісні виміри проєкту.

\subsection{Реєстр ризиків}

Кожен ризик оцінено за двома параметрами: ймовірність виникнення (P) та вплив на проєкт (I), кожен за шкалою 1-3 (1 - низький, 2 - середній, 3 - високий). Значущість (Score) = P $\times$ I.

\vspace{6pt}
\noindent\textbf{Таблиця 12 --- Реєстр ризиків проєкту}

\begin{longtable}{|>{\raggedright\arraybackslash}p{0.05\textwidth}|>{\raggedright\arraybackslash}p{0.29\textwidth}|>{\centering\arraybackslash}p{0.04\textwidth}|>{\centering\arraybackslash}p{0.04\textwidth}|>{\centering\arraybackslash}p{0.07\textwidth}|>{\centering\arraybackslash}p{0.04\textwidth}|>{\raggedright\arraybackslash}p{0.23\textwidth}|}
\hline
\tabhead{ID} & \tabhead{Ризик} & \tabhead{P} & \tabhead{I} & \tabhead{Score} & \tabhead{К} & \tabhead{Стратегія} \\
\hline
\endfirsthead
\hline
\tabhead{ID} & \tabhead{Ризик} & \tabhead{P} & \tabhead{I} & \tabhead{Score} & \tabhead{К} & \tabhead{Стратегія} \\
\hline
\endhead
R01 & Нестача паралельних діалектних даних для тренування & 3 & 3 & \cellcolor{riskred}9 & \cellcolor{riskred}Ч & Пом'якшення \\
\hline
R02 & Низька якість синтетичних даних LLM (галюцинації, помилки) & 3 & 2 & \cellcolor{riskred}6 & \cellcolor{riskred}Ч & Пом'якшення \\
\hline
R03 & Перенавчання (overfitting) через малий обсяг реальних даних & 2 & 3 & \cellcolor{riskred}6 & \cellcolor{riskred}Ч & Пом'якшення \\
\hline
R04 & Недостатні GPU-ресурси для fine-tuning MamayLM & 2 & 2 & \cellcolor{riskyellow}4 & \cellcolor{riskyellow}Ж & Пом'якшення \\
\hline
R05 & Зміна умов доступу до API (OpenRouter, rate limits) & 2 & 2 & \cellcolor{riskyellow}4 & \cellcolor{riskyellow}Ж & Уникнення \\
\hline
R06 & Помилки у вихідних словниках та корпусах & 2 & 2 & \cellcolor{riskyellow}4 & \cellcolor{riskyellow}Ж & Пом'якшення \\
\hline
R07 & Невідповідність BLEU реальній якості нормалізації & 3 & 1 & \cellcolor{riskyellow}3 & \cellcolor{riskyellow}Ж & Прийняття \\
\hline
R08 & Відсутність носіїв діалекту для людської оцінки & 3 & 1 & \cellcolor{riskyellow}3 & \cellcolor{riskyellow}Ж & Пом'якшення \\
\hline
R09 & Зміна вимог керівника або кафедри на пізніх стадіях & 1 & 3 & \cellcolor{riskyellow}3 & \cellcolor{riskyellow}Ж & Пом'якшення \\
\hline
R10 & Технічні збої GPU-сервера під час тривалого навчання & 2 & 1 & \cellcolor{riskgreen}2 & \cellcolor{riskgreen}З & Прийняття \\
\hline
R11 & Складність інтеграції seq2seq і MamayLM в єдиний пайплайн & 1 & 2 & \cellcolor{riskgreen}2 & \cellcolor{riskgreen}З & Пом'якшення \\
\hline
R12 & Обмеження авторських прав на джерельні тексти корпусу & 1 & 2 & \cellcolor{riskgreen}2 & \cellcolor{riskgreen}З & Уникнення \\
\hline
\end{longtable}

\subsection{Матриця управління ризиками}

Матриця розміщує кожен ризик у двовимірному просторі за осями ймовірності та впливу. Колір визначає пріоритет реагування: ЧЕРВОНИЙ - вирішити негайно; ЖОВТИЙ - розробити план дій; ЗЕЛЕНИЙ - моніторинг.

\medskip

\begin{center}
\begin{tabular}{|c|c|c|c|}
\hline
\multirow{2}{*}{\textbf{Ймовірність (P)}} & \multicolumn{3}{c|}{\textbf{Вплив на проєкт (I)}} \\
\cline{2-4}
 & \textbf{1 - низький} & \textbf{2 - середній} & \textbf{3 - високий} \\
\hline
\textbf{3 - високий} &
  \cellcolor{riskyellow}\begin{minipage}[t][1.5cm]{0.22\textwidth}\centering Score~3\\ R07, R08\end{minipage} &
  \cellcolor{riskred}\begin{minipage}[t][1.5cm]{0.22\textwidth}\centering Score~6\\ R02\end{minipage} &
  \cellcolor{riskred}\begin{minipage}[t][1.5cm]{0.22\textwidth}\centering Score~9\\ R01\end{minipage} \\
\hline
\textbf{2 - середній} &
  \cellcolor{riskgreen}\begin{minipage}[t][1.5cm]{0.22\textwidth}\centering Score~2\\ R10\end{minipage} &
  \cellcolor{riskyellow}\begin{minipage}[t][1.5cm]{0.22\textwidth}\centering Score~4\\ R04, R05, R06\end{minipage} &
  \cellcolor{riskred}\begin{minipage}[t][1.5cm]{0.22\textwidth}\centering Score~6\\ R03\end{minipage} \\
\hline
\textbf{1 - низький} &
  \cellcolor{riskgreen}\begin{minipage}[t][1.5cm]{0.22\textwidth}\centering Score~1\\ ---\end{minipage} &
  \cellcolor{riskgreen}\begin{minipage}[t][1.5cm]{0.22\textwidth}\centering Score~2\\ R11, R12\end{minipage} &
  \cellcolor{riskyellow}\begin{minipage}[t][1.5cm]{0.22\textwidth}\centering Score~3\\ R09\end{minipage} \\
\hline
\end{tabular}
\end{center}

\subsection{Плани реагування на ризики}

\textbf{R01 -- Нестача паралельних діалектних даних (Score~9).}
Головна загроза проєкту. Плани пом'якшення: (а)~автоматична генерація синтетичних пар за допомогою LLM на основі правил та словників (реалізовано для гуцульського, BLEU~$\approx$~51); (б)~парсинг відкритих джерел -- закарпатський словник (7\,469 статей); (в)~аугментація через back-translation. Резервний план: обмежити дослідження до двох варіантів (гуцульський + суржик) при критичній нестачі даних.

\textbf{R02 -- Низька якість синтетичних даних (Score~6).}
Плани пом'якшення: (а)~авто\-фільтрація за BLEU-схожістю, довжиною та наявністю діалектних маркерів; (б)~ручна перевірка вибірки 2-5\% на кожній ітерації; (в)~тестування кількох промптів та моделей.

\textbf{R03 -- Перенавчання моделі (Score~6).}
Плани пом'якшення: (а)~QLoRA 4-bit для регуляризації; (б)~моніторинг validation loss; (в)~early stopping; (г)~аугментація та cross-validation. Резервний план: зменшення кількості епох, збільшення dropout.

\textbf{R04 -- Брак GPU-ресурсів (Score~4).} QLoRA знижує потребу до 12-16~ГБ VRAM. Резервний план: оренда GPU через Vast.ai або Google Colab Pro.

\textbf{R05 -- Зміна умов API (Score~4).} Стратегія уникнення: паралельне тестування локальних моделей (Mistral-7B через llama.cpp); локальне збереження всіх згенерованих даних.

\textbf{R06 -- Помилки у словниках (Score~4).} Ручна верифікація 10\% записів; автоматична перевірка через SpaCy та Ukrainian NLP.

\textbf{R07-R09 (Жовті, Score~3).} Прийняття з додатковими заходами: chrF + ручна оцінка 50 речень (R07); залучення студентів-філологів або учасників ТМШ (R08); щотижневі зустрічі з науковим керівником (R09).

\textbf{R10-R12 (Зелені, Score~2).} Регулярний моніторинг: збереження checkpoints кожні 500 кроків (R10); модульна архітектура пайплайну (R11); використання тільки CC-BY корпусів (R12).

\subsection{Зведений план управління ризиками}

\vspace{6pt}
\noindent\textbf{Таблиця 13 --- Зведений план управління ризиками}

\begin{longtable}{|>{\raggedright\arraybackslash}p{0.05\textwidth}|>{\raggedright\arraybackslash}p{0.18\textwidth}|>{\centering\arraybackslash}p{0.04\textwidth}|>{\centering\arraybackslash}p{0.04\textwidth}|>{\centering\arraybackslash}p{0.08\textwidth}|>{\raggedright\arraybackslash}p{0.18\textwidth}|>{\raggedright\arraybackslash}p{0.14\textwidth}|>{\raggedright\arraybackslash}p{0.10\textwidth}|}
\hline
\tabhead{ID} & \tabhead{Ризик (скор.)} & \tabhead{P} & \tabhead{I} & \tabhead{Score} & \tabhead{Дія} & \tabhead{Відповід.} & \tabhead{Термін} \\
\hline
\endfirsthead
\hline
\tabhead{ID} & \tabhead{Ризик (скор.)} & \tabhead{P} & \tabhead{I} & \tabhead{Score} & \tabhead{Дія} & \tabhead{Відповід.} & \tabhead{Термін} \\
\hline
\endhead
R01 & Нестача діалект. даних & 3 & 3 & \cellcolor{riskred}9 & LLM-генерація + парсинг & Коваль Д.П. & Лют. 2026 \\
\hline
R02 & Низька якість LLM-даних & 3 & 2 & \cellcolor{riskred}6 & Авто\-фільтрація + перевірка & Коваль Д.П. & Постійно \\
\hline
R03 & Перенавчання моделі & 2 & 3 & \cellcolor{riskred}6 & QLoRA + early stopping & Коваль Д.П. & До трен. \\
\hline
R04 & Брак GPU-ресурсів & 2 & 2 & \cellcolor{riskyellow}4 & Vast.ai / Colab Pro & Коваль Д.П. & За потр. \\
\hline
R05 & Зміна умов API & 2 & 2 & \cellcolor{riskyellow}4 & Локальні моделі & Коваль Д.П. & Бер. 2026 \\
\hline
R06 & Помилки у словниках & 2 & 2 & \cellcolor{riskyellow}4 & Верифікація 10\% & Коваль Д.П. & Лют. 2026 \\
\hline
R07 & BLEU $\neq$ реальна якість & 3 & 1 & \cellcolor{riskyellow}3 & chrF + human eval & Коваль Д.П. & Кв. 2026 \\
\hline
R08 & Відсутність носіїв мови & 3 & 1 & \cellcolor{riskyellow}3 & Залучення філологів & Коваль Д.П. & Кв. 2026 \\
\hline
R09 & Зміна вимог замовника & 1 & 3 & \cellcolor{riskyellow}3 & Щотижневі зустрічі & Коваль Д.П., Сивоконь О.О. & Постійно \\
\hline
R10 & Збої GPU-сервера & 2 & 1 & \cellcolor{riskgreen}2 & Checkpoints 500 кроків & Коваль Д.П. & Постійно \\
\hline
R11 & Складність інтеграції & 1 & 2 & \cellcolor{riskgreen}2 & Модульна архітектура & Коваль Д.П. & Тр. 2026 \\
\hline
R12 & Авторські права & 1 & 2 & \cellcolor{riskgreen}2 & Тільки CC-BY корпуси & Коваль Д.П. & Лют. 2026 \\
\hline
\end{longtable}

\clearpage

% ==================== РОЗДІЛ 7 ====================
\section{ПОСЛІДОВНІСТЬ ВИКОНАННЯ РОБІТ}

\subsection{Технічний стек та інструменти}

Реалізація проєкту базується на таких технологіях та інструментах:

\begin{itemize}
  \item \textbf{Мова програмування:} Python 3.11+.
  \item \textbf{Фреймворки глибокого навчання:} PyTorch~2.6, HuggingFace Transformers~\cite{hf_transformers}.
  \item \textbf{Parameter-efficient fine-tuning:} PEFT (LoRA / QLoRA)~\cite{lora,qlora}, TRL (SFTTrainer), Unsloth (прискорене навчання для Gemma-3).
  \item \textbf{Управління даними:} HuggingFace Datasets, CSV формат, стратифіковані розбиття.
  \item \textbf{Генерація синтетичних даних:} Mistral Small 24B через OpenRouter API; для суржику додатково rutorphy2 (правило-базована транслітерація).
  \item \textbf{Оцінювання:} SacreBLEU, chrF++, TER; зовнішні LLM (GPT-4o, Gemini 2.0 Flash, Mistral Large) через OpenRouter API.
  \item \textbf{Розпізнавання мовлення:} SpeechBrain (m-ctc-t-large) для транскрипції аудіо.
  \item \textbf{Моніторинг:} TensorBoard, TensorBoardX.
  \item \textbf{Веб-інтерфейс:} NiceGUI~1.x (застосунок <<Лексикон>>).
  \item \textbf{Автоматизація:} GNU Make (Makefile із визначеними цілями).
  \item \textbf{Контроль версій:} Git, GitHub.
\end{itemize}

\subsection{Підготовка даних}

Підготовка даних є критичним (найтривалішим) етапом проєкту (операція B, 6 тижнів).

\textbf{Ручні дані.} Для гуцульського діалекту зібрано понад 10\,000 пар вручну анотованих паралельних прикладів (діалект $\leftrightarrow$ літературна норма). Ці дані використовуються виключно у validation і test розбиттях для об'єктивного оцінювання, а також беруть участь у навчанні.

\textbf{Синтетична генерація.} Основний обсяг навчальних даних отримано через синтетичну генерацію (~30\,000 пар на діалект). Модель \textit{Mistral Small 24B} (через OpenRouter API або локально з 4-bit квантизацією) генерує паралельні пари на основі набору лінгвістичних правил і словника для кожного діалекту. Для суржику додатково застосовується інструмент \textit{rutorphy2} -- правило-базова транслітерація кириличного тексту, що відтворює типові суржикові фонетичні та граматичні явища.

\textbf{Підготовка мульти\-діалектного корпусу.} Скрипт \texttt{prepare\_multidialect\_data.py} завантажує корпуси для всіх чотирьох варіантів, додає префікси-інструкції (<<Переклади з гуцульської:>>, <<Переклади з бойківської:>> тощо), формує стратифіковані розбиття: \textbf{train}~-- 125\,000 пар, \textbf{val}~-- 13\,000 пар, \textbf{test}~-- по 50 вручну верифікованих прикладів на діалект.

\subsection{Навчання моделей}

\textbf{Baseline: umt5-base (seq2seq).} Модель \textit{google/umt5-base} навчається на мульти\-діалектному паралельному корпусі за допомогою скрипту \texttt{train\_encoder\_decoder.py}. Гіперпараметри: 40 епох, batch size 4, gradient accumulation 4, weight decay 0.1, label smoothing 0.1, AdamW 8-bit (BF16), lr~= 5e-5, максимальна довжина послідовності 256 токенів. Навчання виконується на одному GPU без квантизації.

\textbf{Fine-tuning: MamayLM (decoder-only, QLoRA).} Модель \textit{INSAIT-Institute/MamayLM-Gemma-3-4B-IT-v1.0} донавчається за допомогою скрипту \texttt{train\_decoder\_only.py}. Застосовується 4-bit NF4 квантизація (bitsandbytes), QLoRA-адаптери (rank~4, $\alpha$~=~4), SFTTrainer з форматуванням у chat template. Гіперпараметри: 3 епохи, batch size 1, gradient accumulation 2, lr~= 2e-4, max\_seq\_length~= 512. Навчання виконується на GPU через TRL та бібліотеку Unsloth для прискорення.

Прогрес навчання відстежується через TensorBoard: криві loss (training і validation) та метрики BLEU по епохах записуються до директорії \texttt{results/}.

\subsection{Оцінювання та веб-інтерфейс}

\textbf{Автоматичне оцінювання.} Після навчання скрипти \texttt{evaluate\_decoder\_only.py} та \texttt{evaluate\_encoder\_decoder.py} генерують гіпотези на тестових наборах і обчислюють BLEU, chrF++, TER та COMET по кожному діалекту і суржику окремо.

\textbf{Зовнішні моделі.} Для порівняння оцінюються також GPT-4o, Gemini та Mistral-Small через \texttt{eval-dialect} CLI без донавчання (zero-shot / few-shot).

\textbf{Веб-інтерфейс <<Лексикон>>.} Файл \texttt{app.py} реалізує демо-застосунок <<Лексикон>> на NiceGUI. Завантажується основна модель (umt5-base або MamayLM QLoRA). Користувач обирає тип вхідного тексту (Гуцульський / Бойківський / Закарпатський / Суржик), вводить текст або завантажує аудіо -- система повертає нормалізований варіант. Транскрипція аудіо виконується моделлю SpeechBrain (M-CTC-T-large) безпосередньо у браузері.

\subsection{Закриття проєкту та представлення замовникам}

Закриття проєкту передбачає: (1)~завершення написання тексту БКР; (2)~публікацію фінальних моделей на HuggingFace Hub; (3)~публікацію датасету та коду у GitHub-репозиторії з ліцензіями Apache~2.0 і CC-BY~4.0; (4)~підготовку захисної презентації (Idea Pitch, 3-5 хв); (5)~захист бакалаврської роботи на засіданні ДЕК.

\clearpage

% ==================== РОЗДІЛ 8 ====================
\section{ПРЕДСТАВЛЕННЯ ТА СИСТЕМАТИЗАЦІЯ РЕЗУЛЬТАТІВ}

\subsection{Результати навчання моделей}

Нижче наведено зведену таблицю фактично досягнутих метрик якості для всіх підтримуваних варіантів вхідного тексту. Порівнюються три підходи: fine-tuned umt5-base (seq2seq), MamayLM з QLoRA (decoder-only) та GPT-4o Mini у режимі zero-shot (без донавчання). Нижче TER -- краще; вище BLEU і chrF++ -- краще.

\vspace{6pt}
\noindent\textbf{Таблиця 14 --- Метрики якості нормалізації (тестовий набір, 50 прикладів/діалект)}

\begin{longtable}{|>{\raggedright\arraybackslash}p{0.22\textwidth}|>{\raggedright\arraybackslash}p{0.22\textwidth}|>{\centering\arraybackslash}p{0.12\textwidth}|>{\centering\arraybackslash}p{0.12\textwidth}|>{\centering\arraybackslash}p{0.12\textwidth}|}
\hline
\tabhead{Діалект} & \tabhead{Модель} & \tabhead{BLEU$\uparrow$} & \tabhead{chrF++$\uparrow$} & \tabhead{TER$\downarrow$} \\
\hline
\endfirsthead
\hline
\tabhead{Діалект} & \tabhead{Модель} & \tabhead{BLEU$\uparrow$} & \tabhead{chrF++$\uparrow$} & \tabhead{TER$\downarrow$} \\
\hline
\endhead
\multirow{3}{*}{Гуцульський} & umt5-base & \textbf{74.68} & \textbf{86.90} & \textbf{12.53} \\
\cline{2-5}
 & MamayLM QLoRA & 64.84 & 74.02 & 37.60 \\
\cline{2-5}
 & GPT-4o (zero-shot) & 64.89 & -- & -- \\
\hline
\multirow{2}{*}{Бойківський} & umt5-base & 29.12 & 54.78 & 49.37 \\
\cline{2-5}
 & MamayLM QLoRA & \textbf{51.81} & \textbf{66.82} & \textbf{35.26} \\
\hline
\multirow{3}{*}{Закарпатський} & umt5-base & \textbf{6.20} & \textbf{27.45} & \textbf{82.21} \\
\cline{2-5}
 & MamayLM QLoRA & 3.82 & 20.93 & 89.69 \\
\cline{2-5}
 & GPT-4o (zero-shot) & 65.39 & -- & -- \\
\hline
\multirow{3}{*}{Суржик} & umt5-base & \textbf{89.42} & \textbf{93.58} & \textbf{5.98} \\
\cline{2-5}
 & MamayLM QLoRA & 78.48 & 91.22 & 11.96 \\
\cline{2-5}
 & GPT-4o (zero-shot) & 88.97 & -- & -- \\
\hline
\end{longtable}

\noindent \textit{Примітка: жирним виділено найкращий результат по кожному діалекту. <<-->> означає, що показник не обчислювався для даної моделі. GPT-4o та Gemini оцінювались лише за BLEU через обмеження платного API.}

\subsection{Загальні висновки за результатами моделювання}

Результати експериментів (Таблиця~14) демонструють суттєву різницю у складності нормалізації між діалектами. \textbf{Суржик} виявився найлегшим для обох моделей: umt5-base досяг BLEU~=~89.42, що зіставно з рівнем GPT-4o (88.97) без будь-якого донавчання. \textbf{Гуцульський діалект} також показав сильні результати у seq2seq підходу (BLEU~=~74.68, TER~=~12.53), тоді як MamayLM після донавчання поступається на 10 BLEU-балів -- ймовірно, через особливості chat-template формату. \textbf{Бойківський діалект} виявив зворотну картину: MamayLM~QLoRA перевершив umt5-base на 22.7~BLEU (51.81 проти 29.12), що вказує на більшу здатність decoder-only архітектури до узагальнення для менш ресурсних варіантів. \textbf{Закарпатський діалект} залишається найскладнішим для fine-tuned моделей (BLEU~$<$~7), але GPT-4o нульовим пострілом досягнув BLEU~=~65.39, що свідчить про нестачу ресурсів у тренувальному корпусі та вказує напрямок для майбутніх досліджень.

Ключовими факторами якості є: (1)~стратифіковане балансування корпусу за діалектами; (2)~синтетична генерація через Mistral Small 24B за правило-базованими промптами; (3)~використання лише вручну верифікованих даних у validation і test розбиттях. Застосунок <<Лексикон>> забезпечує end-to-end демонстрацію від голосового введення до нормалізованого тексту.

\clearpage

% ==================== ВИСНОВКИ ====================
\section*{ВИСНОВКИ}
\addcontentsline{toc}{section}{ВИСНОВКИ}

У ході виконання розрахунково-графічної роботи сформовано повну технічну та комунікаційну документацію для IT-проєкту <<Система для автоматичної нормалізації суржику та діалектних форм української мови>>. Нижче підведено підсумки за всіма 13 пунктами змістового контенту.

\begin{enumerate}
  \item \textbf{Вступ.} Обгрунтовано актуальність задачі нормалізації діалектів і суржику в контексті сучасного стану українського суспільства та розвитку NLP-технологій.

  \item \textbf{Сфера застосування та об'єм робіт.} За методами «Критичних запитань» і SMART визначено: проблему, потреби, сферу, зацікавлені сторони, межі проєкту та трикутник обмежень (час / ресурси / якість).

  \item \textbf{Процес управління та концептуальні поняття.} Проєкт визначено як тимчасове зусилля зі створення унікального програмного продукту. Описано команду (виконавець + науковий керівник) і ромб конкурентних обмежень (кошти, час, якість, результат).

  \item \textbf{Складність проєкту.} Проєкт класифіковано як еволюційний із середнім рівнем виклику та менеджментною критичністю.

  \item \textbf{Вимоги до проєкту (RAM).} Сформовано матрицю розподілу відповідальності для 13 ключових завдань і 10 стейкхолдерів.

  \item \textbf{Стейкхолдери.} Ідентифіковано 10 зацікавлених сторін -- від виконавця і наукового керівника до носіїв діалектів, NLP-спільноти та EdTech-платформ. Побудовано сітку Power/Interest.

  \item \textbf{Вимоги стейкхолдерів та планування.} Сформовано таблицю вимог, прив'язану до фаз. Для кожної з п'яти фаз PMBOK розроблено контрольний список.

  \item \textbf{Контрольний список та WBS.} Побудовано структуру розбиття робіт із 6 пакетів і 13 завдань. Реалізовано 16-крокову перевірку планування.

  \item \textbf{Аналіз ризиків.} Ідентифіковано 12 ризиків, оцінено ймовірність та вплив, побудовано матрицю і розроблено конкретні плани реагування.

  \item \textbf{Комунікаційна матриця.} Визначено аудиторії, типи повідомлень, частоту та відповідальних для 6 каналів комунікації.

  \item \textbf{Закриття проєкту.} Визначено артефакти закриття: публікація БКР, відкритий репозиторій, модель на HuggingFace, захист.

  \item \textbf{Представлення проєкту.} Підготовлено Idea Pitch (3-5 хв) за структурою: проблема, рішення, результати, висновки.

  \item \textbf{Підсумки.} Навчання завершено для обох моделей на всіх чотирьох діалектах. umt5-base досяг BLEU~=~89.42 на суржику та 74.68 на гуцульському. MamayLM~QLoRA перевершив baseline на бойківському (BLEU~51.81 проти 29.12). Закарпатський діалект залишається викликом для донавчених моделей, але GPT-4o (zero-shot) демонструє BLEU~=~65.39, що вказує на перспективи збільшення корпусу. Застосунок <<Лексикон>> розгорнуто та готовий до демонстрації.
\end{enumerate}

\clearpage

% ==================== СПИСОК ДЖЕРЕЛ ====================
\section*{СПИСОК ВИКОРИСТАНИХ ДЖЕРЕЛ}
\addcontentsline{toc}{section}{СПИСОК ВИКОРИСТАНИХ ДЖЕРЕЛ}

\begin{thebibliography}{99}

\bibitem{pmbok}
Project Management Institute.
\textit{A Guide to the Project Management Body of Knowledge (PMBOK Guide)}.
7th~ed.~--- PMI, 2021.~--- 370~p.

\bibitem{vaswani}
Vaswani~A., Shazeer~N., Parmar~N. et al.
Attention is All You Need~//
\textit{Advances in Neural Information Processing Systems 30 (NeurIPS~2017)}.~---
Curran Associates, 2017.~--- P.~5998--6008.

\bibitem{raffel}
Raffel~C., Shazeer~N., Roberts~A. et al.
Exploring the Limits of Transfer Learning with a Unified Text-to-Text Transformer~//
\textit{Journal of Machine Learning Research}.~--- 2020.~--- Vol.~21(140).~--- P.~1--67.

\bibitem{lora}
Hu~E.J., Shen~Y., Wallis~P. et al.
LoRA: Low-Rank Adaptation of Large Language Models~//
\textit{Proceedings of the International Conference on Learning Representations (ICLR~2022)}.~---
OpenReview, 2022.

\bibitem{qlora}
Dettmers~T., Pagnoni~A., Holtzman~A., Zettlemoyer~L.
QLoRA: Efficient Finetuning of Quantized LLMs~//
\textit{Advances in Neural Information Processing Systems 36 (NeurIPS~2023)}.~---
Curran Associates, 2023.

\bibitem{uagec}
Syvokon~O., Nahorna~O.
UA-GEC: Grammatical Error Correction and Fluency Corpus for the Ukrainian Language~//
\textit{Proceedings of the Second Ukrainian Natural Language Processing Workshop (UNLP~2023)}.~---
Association for Computational Linguistics, 2023.~--- P.~143--149.

\bibitem{hf_transformers}
Wolf~T., Debut~L., Sanh~V. et al.
Transformers: State-of-the-Art Natural Language Processing~//
\textit{Proceedings of the 2020 Conference on Empirical Methods in Natural Language Processing (EMNLP~2020): System Demonstrations}.~---
Association for Computational Linguistics, 2020.~--- P.~38--45.

\end{thebibliography}

\end{document}
